\renewcommand{\thesection}{S5}
\section{Material Selection}
\setcounter{subsection}{4}
\subsection{Ranking}
    From the previously defined constraints (presented in Table 6 of the main report), 59 materials from the databases satisfied the set restrictions. All these materials were part of the metals family and included aluminium alloys, titanium alloys, nickel-titanium alloys, nickel-chromium alloys, chromium metal, zirconium alloys, among others. The evaluation focused on yield strength and density to rank the materials effectively.
    
    A performance index, shown in equation \ref{eq:performance}, was derived from equation 1 of the main report to refine the analysis further and applied to generate an Ashby plot for the previously mentioned materials, this plot is represented in Figure 16 of the main report. This plot facilitated visualization and comparison of material performance based on the performance index.

    \renewcommand{\theequation}{S2}
    \begin{equation}
        \log(\sigma) = \frac{3}{2}\left(\log(\rho) + \log(P_{\text{cr}})\right)
        \label{eq:performance}
    \end{equation}

    This equation allowed for the plotting of lines with a constant slope which served as a tool to ascertain the materials with the highest $P_{cr}$ (the highest the value of the line at $x=0$, the highest the $P_{cr}$).